\chapter{Mots-clés VBS / CATScript}

\noindent Référence des objets et méthodes principaux utilisés dans les macros CATIA~V5 (CATScript / VBScript).

\vspace{1.5em}

% -------------------------------------------------------
% Table 1 — Objets principaux
% -------------------------------------------------------
\begin{longtable}{>{\ttfamily}p{5.5cm} p{9.5cm}}
  \caption{Objets principaux}\label{annex:catia-1} \\
  \hline
  \normalfont\textbf{Mot-clé} & \textbf{Description} \\
  \hline
  \endfirsthead
  \multicolumn{2}{r}{\small\textit{(suite\dots)}} \\
  \hline
  \normalfont\textbf{Mot-clé} & \textbf{Description} \\
  \hline
  \endhead
  \hline
  \endlastfoot
  %
  CATIA          & Objet racine de l'application CATIA. Point d'entrée de toute macro --- donne accès aux documents, paramètres et sessions. \\[3pt]
  ActiveDocument & Retourne le document actuellement ouvert et actif dans CATIA (Part, Product, Drawing, etc.). \\[3pt]
  PartDocument   & Type de document représentant un fichier \normalfont\texttt{.CATPart}. Permet l'accès à la géométrie solide et aux features. \\[3pt]
  Part           & Objet principal d'un PartDocument. Contient le PartBody, les paramètres, les relations et la géométrie. \\[3pt]
  MainBody       & Corps solide principal d'un Part (\normalfont\texttt{PartBody}). C'est là où sont stockées toutes les features solides. \\[3pt]
  Body           & Objet représentant un corps solide (PartBody ou corps secondaire). Contient une collection de shapes/features. \\[3pt]
  Shapes         & Collection de toutes les features solides (Pad, Shaft, Pocket\dots) contenues dans un Body. \\[3pt]
  HybridBodies   & Collection des corps hybrides (contenant surfaces + wireframe) dans un Part ou un Body. \\[3pt]
  HybridShapes   & Collection des éléments surfaciques et filaires (splines, extrusions, fills\dots) dans un HybridBody. \\
\end{longtable}

\vspace{1em}

% -------------------------------------------------------
% Table 2 — Mesure & Analyse (SPA)
% -------------------------------------------------------
\begin{longtable}{>{\ttfamily}p{5.5cm} p{9.5cm}}
  \caption{Mesure \& Analyse (SPA)}\label{annex:catia-2} \\
  \hline
  \normalfont\textbf{Mot-clé} & \textbf{Description} \\
  \hline
  \endfirsthead
  \multicolumn{2}{r}{\small\textit{(suite\dots)}} \\
  \hline
  \normalfont\textbf{Mot-clé} & \textbf{Description} \\
  \hline
  \endhead
  \hline
  \endlastfoot
  %
  SPAWorkbench    & Workbench \textit{Space Analysis} de CATIA. Fournit les outils de mesure géométrique (volume, aire, centre de masse). \\[3pt]
  GetWorkbench()  & Méthode du document pour accéder à un workbench spécifique. Ex : \normalfont\texttt{oDoc.GetWorkbench("SPAWorkbench")}. \\[3pt]
  GetMeasurable() & Méthode du SPAWorkbench qui crée un objet \normalfont\texttt{Measurable} à partir d'un Body ou d'une surface. \\[3pt]
  Measurable      & Objet retourné par \normalfont\texttt{GetMeasurable()}. Expose les propriétés de mesure : volume, aire, centre de masse. \\[3pt]
  .Volume         & Propriété de \normalfont\texttt{Measurable}. Retourne le volume du solide en m$^{3}$ (à convertir en mm$^{3}$ via $\times 10^{9}$). \\[3pt]
  .Area           & Propriété de \normalfont\texttt{Measurable}. Retourne l'aire totale de la surface du solide en m$^{2}$ ($\times 10^{6}$ pour mm$^{2}$). \\[3pt]
  GetCOG()        & Méthode de \normalfont\texttt{Measurable}. Retourne les coordonnées du centre de gravité (Center Of Gravity) du solide. \\
\end{longtable}

\vspace{1em}

% -------------------------------------------------------
% Table 3 — Surfaces & Géométrie (GSD)
% -------------------------------------------------------
\begin{longtable}{>{\ttfamily}p{5.5cm} p{9.5cm}}
  \caption{Surfaces \& Géométrie (GSD)}\label{annex:catia-3} \\
  \hline
  \normalfont\textbf{Mot-clé} & \textbf{Description} \\
  \hline
  \endfirsthead
  \multicolumn{2}{r}{\small\textit{(suite\dots)}} \\
  \hline
  \normalfont\textbf{Mot-clé} & \textbf{Description} \\
  \hline
  \endhead
  \hline
  \endlastfoot
  %
  GSMWorkbench              & Workbench \textit{Generative Shape Design}. Donne accès aux outils de création et d'analyse surfacique. \\[3pt]
  HybridShapeFactory        & Fabrique d'objets surfaciques dans GSD. Permet de créer des surfaces, courbes et points par macro. \\[3pt]
  Reference                 & Objet qui encapsule une entité géométrique (face, arête, surface) pour l'utiliser dans les méthodes CATIA. \\[3pt]
  CreateReferenceFromObject() & Crée un objet \normalfont\texttt{Reference} à partir d'un objet géométrique, nécessaire pour beaucoup d'opérations GSD. \\[3pt]
  SurfaceType               & Propriété indiquant le type d'une surface (plane, cylindrique, NURBS\dots). Utile pour classer les géométries. \\[3pt]
  Face                      & Représente une face individuelle d'un solide ou d'une surface. Peut être analysée séparément. \\
\end{longtable}

\vspace{1em}

% -------------------------------------------------------
% Table 4 — Contrôle de flux VBS
% -------------------------------------------------------
\begin{longtable}{>{\ttfamily}p{5.5cm} p{9.5cm}}
  \caption{Contrôle de flux VBS}\label{annex:catia-4} \\
  \hline
  \normalfont\textbf{Mot-clé} & \textbf{Description} \\
  \hline
  \endfirsthead
  \multicolumn{2}{r}{\small\textit{(suite\dots)}} \\
  \hline
  \normalfont\textbf{Mot-clé} & \textbf{Description} \\
  \hline
  \endhead
  \hline
  \endlastfoot
  %
  TypeName()    & Fonction VBS qui retourne le type d'un objet sous forme de chaîne. Utilisé pour vérifier si un document est un \normalfont\texttt{PartDocument}. \\[3pt]
  On Error GoTo & Directive de gestion d'erreurs. Redirige l'exécution vers un label défini si une erreur survient. \\[3pt]
  MsgBox        & Affiche une boîte de dialogue à l'utilisateur avec un message, une icône et un titre personnalisables. \\[3pt]
  Set           & Mot-clé VBS pour affecter un objet à une variable. Obligatoire pour tous les objets CATIA (pas les types simples). \\[3pt]
  Nothing       & Valeur nulle pour les objets VBS. Utilisé en fin de macro pour libérer les références mémoire (\normalfont\texttt{Set obj = Nothing}). \\[3pt]
  Format()      & Formate un nombre en chaîne de caractères avec un masque. Ex : \normalfont\texttt{Format(val, "\#,\#\#0.000")} pour l'affichage. \\[3pt]
  Chr(13)       & Caractère retour à la ligne dans une chaîne VBS. Utilisé pour structurer les messages \normalfont\texttt{MsgBox}. \\
\end{longtable}

\vspace{1em}

% -------------------------------------------------------
% Table 5 — Sélection & Interaction
% -------------------------------------------------------
\begin{longtable}{>{\ttfamily}p{5.5cm} p{9.5cm}}
  \caption{Sélection \& Interaction}\label{annex:catia-5} \\
  \hline
  \normalfont\textbf{Mot-clé} & \textbf{Description} \\
  \hline
  \endfirsthead
  \multicolumn{2}{r}{\small\textit{(suite\dots)}} \\
  \hline
  \normalfont\textbf{Mot-clé} & \textbf{Description} \\
  \hline
  \endhead
  \hline
  \endlastfoot
  %
  Selection        & Objet CATIA représentant la sélection active dans l'interface. Permet de récupérer ou définir des éléments sélectionnés. \\[3pt]
  SelectElement2() & Méthode de \normalfont\texttt{Selection} qui invite l'utilisateur à sélectionner un élément dans l'interface CATIA via un filtre. \\[3pt]
  Item(i)          & Méthode générique pour accéder au $i$-ème élément d'une collection (Shapes, Bodies, HybridBodies\dots). \\[3pt]
  .Count           & Propriété de toute collection CATIA. Retourne le nombre d'éléments qu'elle contient. \\[3pt]
  .Name            & Propriété universelle. Retourne ou définit le nom d'un objet CATIA (Part, Body, feature, etc.). \\
\end{longtable}